\chapter{Proofs for Heuristic Transformations and Working Backwards}
\label{annex:p12-heuristic-transformations-and-backward-reasoning}

\begin{annexitem}{annex:p12.heuristic.solution-relations}{thm:p12-heuristic-solution-relations}{Solution sets and reconstruction maps}
\begin{proof}
If \(x\in S_C\), then \(C(x)\), so \(C(x)\Rightarrow N(x)\)
gives \(x\in S_N\). This proves the first inclusion. Applying
\(H(x)\Rightarrow C(x)\) to a member of \(S_H\) gives the
second. The two directions of an equivalence give both inclusions
and hence equality. These statements use the same admissible domain;
an altered domain must be translated as well as its predicate.

By the truth condition of finite conjunction,
\(x\in S_C\) holds exactly when \(x\) belongs to every
\(S_{C_i}\), proving the intersection identity. For \(k=0\),
the condition is true and the empty intersection relative to the
specified domain is that domain. On \(X=\{0,1\}\), the
conditions \(x=0\) and \(x=1\) have nonempty solution sets
but empty intersection, so separately solved parts need not combine.

For subproblems on coordinate sets \(I_1,\ldots,I_k\), take
their solution tuples \(x_i\). If the tuples agree on every
overlap, define a tuple on \(\bigcup_i I_i\) by choosing the
value assigned by any containing tuple. Agreement makes this
well defined. Its restriction to \(I_i\) is \(x_i\), so it
satisfies each local condition. It satisfies the original problem
only if all additional linking conditions hold too. Conversely,
restriction of a global solution gives compatible local tuples.
This states precisely what recombination of coordinate subproblems
requires; it does not assume that arbitrary local witnesses agree.

For the auxiliary map, the assumption \(C_d(F(y))\) for
\(y\in A\) gives \(F(A)\subseteq S_d\). Thus each actual
auxiliary solution yields an original solution. If \(G\) has
the stated right-inverse property, then every \(x\in S_d\)
equals \(F(G(x))\) and belongs to \(F(A)\). The reverse
inclusion already proved gives equality. Without surjectivity
one may find some original solutions but has not thereby enumerated all.
\end{proof}
\end{annexitem}

\begin{annexitem}{annex:p12.heuristic.triangle-loci}{thm:p12-heuristic-triangle-loci}{The two geometric conditions and all solution cases}
\begin{proof}
The prescribed upper-half-plane distance from the line \(BC\)
forces \(A=(t,h)\) for some real \(t\). The vectors from
\(A\) to \(B,C\) are \((-t,-h)\) and \((a-t,-h)\).
Their dot product is \(t^2-at+h^2\), and their determinant
is the positive number \(ah\). Consequently, if \(\theta\)
is the interior angle at \(A\), then
\[
              \cot\theta=\frac{t^2-at+h^2}{ah}.
\]
Indeed cosine is the dot product divided by the product of the
vector lengths, while sine is the positive determinant divided
by that same product. Because \(a,h>0\), neither vector vanishes
and \(0<\theta<\pi\). Cotangent is one-to-one on this interval.
Thus \(\theta=\alpha\) exactly when
\[
 t^2-at+h^2-ah\cot\alpha=0,
 \quad\text{or}\quad
 (t-a/2)^2=a^2/4-h^2+ah\cot\alpha.
\]
The three signs of the right side give zero, one or two real
solutions as stated. Every resulting point has the required
altitude and angle, so the argument both generates and verifies
all candidates. Reflection across the base gives the other
half-plane without altering lengths or the interior angle.

For a varying upper point \((t,s)\), the angle equation is
\[
 (t-a/2)^2+(s-(a/2)\cot\alpha)^2
                =\frac{a^2}{4\sin^2\alpha},\qquad s>0.
\]
It is the relevant open arc of a circle through \(B,C\), with
the endpoints excluded. A straightedge-and-compass description
can first construct an isosceles triangle on \(BC\) with base
angles \((\pi-\alpha)/2\), then its circumcircle using two
perpendicular bisectors. The portion on the chosen side with
the required angle is the locus just calculated. Intersect it
with the parallel at height \(h\). Choosing an arbitrary point
on each locus separately would not impose both requirements on
one vertex.
\end{proof}
\end{annexitem}

\begin{annexitem}{annex:p12.heuristic.box}{thm:p12-heuristic-box}{Two right triangles and the original length conditions}
\begin{proof}
Place the box vertices at the appropriate combinations of
\(0,a\), \(0,b\), and \(0,c\) on three mutually perpendicular
coordinate axes. The face diagonal from \((0,0,0)\) to
\((a,b,0)\) has squared length \(a^2+b^2\) by Pythagoras.
Its vector \((a,b,0)\) is perpendicular to \((0,0,c)\):
their dot product is zero. They therefore form the legs of a
right triangle with the space diagonal as hypotenuse. A second
application gives \(x^2=a^2+b^2+c^2\). A length is positive,
so only the positive square root is admissible.

Opposite vertices of each space diagonal have midpoint
\((a/2,b/2,c/2)\), proving the common center and half-length
assertion. Permuting the dimensions does not alter their sum
of squares. If a positive dimension is increased, its square
strictly increases, as does the nonnegative square root of the
sum. Multiplication of all dimensions by \(\lambda>0\) multiplies
the squared sum by \(\lambda^2\) and its positive root by
\(\lambda\). Continuity of square root on positive numbers
gives the planar limit when \(c\) tends to zero. Each check
follows from the already derived expression and its domain.

For the roof, the horizontal displacement from center to a corner
has component lengths \(21/2=10.5\) and \(16/2=8\). The cable
attachment is \(8-2=6\) yards above the roof. Hence its squared
length is
\(10.5^2+8^2+6^2=110.25+64+36=210.25=14.5^2\).
Each corner changes only the signs of the two horizontal
components, so all four cable lengths agree. This verifies the
stated geometric model without adding a practical allowance to it.
\end{proof}
\end{annexitem}

\begin{annexitem}{annex:p12.heuristic.geometric-specifications}{thm:p12-heuristic-geometric-specifications}{Side constraints, altitude constraints and a proof target}
\begin{proof}
In a nondegenerate triangle each side is strictly shorter than
the sum of the other two. The proposed side length \(3\) is
greater than \(1+1\), so that side specification is impossible.
For the stated coordinate triangle the base length is
\(6/\sqrt{35}\), the altitude to it is \(3\), and its area
is \(9/\sqrt{35}\). Each other side has length
\[
       \sqrt{9+9/35}=18/\sqrt{35}.
\]
The altitude to either of those sides is twice the area divided
by that side length, hence
\((18/\sqrt{35})/(18/\sqrt{35})=1\). All vertices are distinct
and the area is positive, so this is an admissible triangle with
the required three altitudes.

For the quadrilateral, name its consecutive vertices \(A,B,C,D\)
and their position vectors \(a,b,c,d\). Equal adjacent side
lengths at \(B\) imply
\(\|b-a\|^2=\|b-c\|^2\), and expansion yields
\(2b\cdot(c-a)=\|c\|^2-\|a\|^2\).
The equal side lengths at \(D\) give the same equation with
\(d\) replacing \(b\). Subtraction gives
\((b-d)\cdot(c-a)=0\). Since both diagonal vectors are
nonzero, this is their perpendicularity. The argument identifies
exactly where the side-equality hypothesis is used.

For the final triangle equivalence, the cosine law follows by
expanding the squared norm of the difference of its two side
vectors. If all sides are equal, it gives cosine \(1/2\)
for each interior angle, so each is \(\pi/3\). Conversely,
three equal interior angles each equal \(\pi/3\), since their
sum is \(\pi\). For opposite side lengths \(a,b,c\), the
cosine law gives \(a^2=b^2+c^2-bc\) and
\(b^2=a^2+c^2-ac\). Subtracting yields
\((a-b)(2(a+b)-c)=0\). The triangle inequality makes the
second factor positive, so \(a=b\). Repeating the argument
with another pair gives \(a=b=c\), proving the equivalence.
\end{proof}
\end{annexitem}

\begin{annexitem}{annex:p12.heuristic.backward-proof}{thm:p12-heuristic-backward-proof}{Forward justification of a backwards-discovered proof}
\begin{proof}
Start with the derivation of \(T_n\) from \(\Gamma\).
Modus ponens with \(T_n\to T_{n-1}\) gives \(T_{n-1}\).
Applying the remaining supplied implications successively gives
\(T_{n-2},\ldots,T_0\), using the same assumptions throughout.
For \(n=0\), the supplied derivation already is the desired one.
Each equivalence supplies both implications, so a chain of
equivalences can also be traversed backwards under its side conditions.

For the converse-direction failure, consider the target \(x=1\)
and its necessary consequence \(x^2=1\) over the reals. The value
\(x=-1\) satisfies the latter and not the former. Thus obtaining
the consequence, even from correct reasoning, does not establish
the target. A temporary assumption of the target can help discover
necessary relations; it must not remain as an undischarged premise
of the supposed final proof.
\end{proof}
\end{annexitem}

\begin{annexitem}{annex:p12.heuristic.square-root}{thm:p12-heuristic-square-root}{The retained sign guard and final substitution}
\begin{proof}
The real square root requires \(x+1\ge0\), and its value is
nonnegative. Equality to \(x-1\) therefore requires \(x\ge1\),
which also implies the square root's domain condition. Squaring
a valid equality gives
\[
 x+1=(x-1)^2=x^2-2x+1,
 \qquad x^2-3x=x(x-3)=0.
\]
Since a product of real numbers is zero only if one factor is
zero, its candidates are \(0\) and \(3\). The first violates
\(x\ge1\) and directly gives \(1=-1\) in the original equation.
The second gives \(\sqrt4=2\) and is valid.

For reversibility with the guard, if \(x\ge1\) and
\(x+1=(x-1)^2\), taking the nonnegative square root gives
\(\sqrt{x+1}=|x-1|=x-1\). Thus the guarded squared equation
is equivalent to the original one. This proves both validity of
the returned answer and absence of any lost solution.
\end{proof}
\end{annexitem}

\begin{annexitem}{annex:p12.heuristic.backward-path}{thm:p12-heuristic-backward-path}{Finite paths and predecessor witnesses}
\begin{proof}
For \(k=0\), membership in \(G_0=G\) is exactly a path of
length zero to a goal. Suppose the claim holds for \(k\).
Membership of \(s\) in \(G_{k+1}\) means there is a \(t\in G_k\)
with \((s,t)\in R\). By induction there is a legal \(k\)-step
path from \(t\) to a goal; prefixing its edge from \(s\) gives
the required \((k+1)\)-step path. Conversely, the first edge
of any such path leads to a state whose remaining \(k\)-step
path places it in \(G_k\), again by induction. Hence the
initial state is in \(\operatorname{Pre}_R(G_k)\).

The proof is constructive for a recorded finite chain of
witnesses: list its states in the order of the certified forward
edges. No assumption that every action has a single-valued inverse
has been used.
\end{proof}
\end{annexitem}

\begin{annexitem}{annex:p12.heuristic.pails}{thm:p12-heuristic-pails}{Every legal action in the pail solution}
\begin{proof}
All amounts in the displayed sequence lie within the two capacities.
Fill the large container to obtain \((9,0)\). Pour four quarts
into the empty small container, yielding \((5,4)\), and empty
the small one, yielding \((5,0)\). Another pour fills the small
container and leaves \((1,4)\); emptying it gives \((1,0)\).
Pour the remaining one quart from the large container into the
small one: this pour ends when its source is empty and gives
\((0,1)\). Fill the large container to obtain \((9,1)\).
The small container can now receive only three quarts, so pouring
until it is full leaves \((6,4)\). Emptying it gives \((6,0)\).
Every action uses a full capacity or an empty-source/full-destination
stopping condition; no intermediate marking or guessed amount is used.

The backward insight is certified by its forward edge:
\((9,1)\to(6,4)\to(6,0)\). The earlier legal actions certify
reachability of \((9,1)\) from the initial state. Concatenating
the two paths proves that the construction meets the original
capacity and operation conditions. No assertion of a shortest
possible action sequence is needed or made.
\end{proof}
\end{annexitem}

\begin{annexitem}{annex:p12.heuristic.generalization}{thm:p12-heuristic-generalization}{Restriction, counterexamples and finite observations}
\begin{proof}
Every \(e\in E\) is in \(D\), so universal truth on \(D\)
gives \(P(e)\). This proves the restriction. If some \(e\in E\)
instead falsifies \(P\), it is also a counterexample in \(D\).
For the failed reverse implication, take \(D=\{0,1\}\),
\(E=\{0\}\), and \(P(x)\) to mean \(x=0\).
The restricted statement is true but the statement on \(D\) is false.

For \(1\le j\le N\), one factor of \(p_N(j)\) is zero,
so the product is zero. At \(N+1\), the factors are
\(N,N-1,\ldots,1\); their product is the positive integer
\(N!\). Thus even an arbitrary finite number of correct sample
values is compatible with a false universal identity. A separate
bound on polynomial degree, for example, would be additional
information and is not part of this example's observations alone.
\end{proof}
\end{annexitem}

\begin{annexitem}{annex:p12.heuristic.cubes}{thm:p12-heuristic-cubes}{Induction for the sum of cubes}
\begin{proof}
For \(n=0\), the empty sum and the displayed square are zero.
Assume the equality at \(n\). Adding \((n+1)^3\) gives
\[
\begin{aligned}
 \sum_{k=1}^{n+1}k^3
 &=\frac{n^2(n+1)^2}{4}+(n+1)^3\\
 &=\frac{(n+1)^2[n^2+4(n+1)]}{4}\\
 &=\frac{(n+1)^2(n+2)^2}{4}
  =\left(\frac{(n+1)(n+2)}2\right)^2.
\end{aligned}
\]
Mathematical induction establishes the formula for all
nonnegative integers. In particular \(n=4\) gives \(10^2=100\).
The observed special case motivates the conjecture; the base case
and the proved implication from \(n\) to \(n+1\) justify it.
\end{proof}
\end{annexitem}

\begin{annexitem}{annex:p12.heuristic.central-bisection}{thm:p12-heuristic-central-bisection}{Central symmetry and the line-through-center case}
\begin{proof}
Write the solid as \(K\) and a plane through \(o\) as
\(\{x:n\cdot(x-o)=0\}\), where \(n\ne0\).
The central inversion \(J(x)=2o-x\) sends \(K\) onto itself
and interchanges the portions in the two open half-spaces,
because \(n\cdot(J(x)-o)=-n\cdot(x-o)\).
It is an isometry and therefore preserves Euclidean volume.
Equivalently, its linear part is \(-I\), whose determinant
has absolute value one. The two portions have equal volume.
The intervening plane has three-dimensional volume zero:
inside any bounded box its portion can be enclosed in slabs
of arbitrarily small thickness. Thus the two portions each
have half the volume of \(K\).

Let \(\ell=\{p+tu:t\in\mathbb R\}\), \(u\ne0\).
If \(o\notin\ell\), the independent vectors \(u,o-p\)
span the unique plane through the line and center. If
\(o\in\ell\), add any direction independent of \(u\)
to obtain such a plane; there are many choices, and the same
volume argument applies to each. A regular octahedron has a
center of symmetry, so it is one instance of this general result.
\end{proof}
\end{annexitem}

\begin{annexitem}{annex:p12.heuristic.frustum}{thm:p12-heuristic-frustum}{A parameterized volume calculation with its division guard}
\begin{proof}
Let the full pyramid above the larger base have perpendicular
height \(H\). Similarity gives the side ratio of its smaller
removed pyramid as \(b/a\), so its height is \(Hb/a\).
The remaining height is \(h=H-Hb/a\). Since \(a>b\ge0\),
\(a-b>0\) and this equation can be solved as
\(H=ha/(a-b)\), with removed height \(hb/(a-b)\).

A pyramid with base area \(A\) and height \(H\) has volume
\(AH/3\): its section at height \(z\) above the base is
similar to the base with scale \(1-z/H\), and hence area
\(A(1-z/H)^2\). Integrating these areas from \(0\) to \(H\)
gives \(AH/3\). This is also the elementary pyramid-volume
formula used in the calculation. Subtracting the two pyramid
volumes now gives
\[
 V=\frac h3\frac{a^3-b^3}{a-b}
  =\frac h3(a^2+ab+b^2),
\]
where factorization of \(a^3-b^3\) and the positive denominator
justify the cancellation. If \(b=0\), the removed pyramid has
zero volume and the same expression applies. If \(a=b>0\), the
division is not allowed: the object is a prism and its base
area times height is \(a^2h\), equal to the asserted formula
by direct substitution. Finally, \((6/3)(100+50+25)=350\).
Scaling all lengths by \(\lambda>0\) scales this expression
by \(\lambda^3\), as required for volume.
\end{proof}
\end{annexitem}

\begin{annexitem}{annex:p12.heuristic.triangle-counterexample}{thm:p12-heuristic-triangle-counterexample}{A nondegenerate counterexample with exact radii}
\begin{proof}
Normalize an equilateral triangle to vertices \((-1,0),(1,0),
(0,\sqrt3)\). Its center on the symmetry axis is
\((0,\sqrt3/3)\): the distances to the three vertices are
\(2\sqrt3/3\), and the distances to its three side lines are
\(\sqrt3/3\), by direct Euclidean distance calculations.
Thus \(r=H/3\) and \(R=2H/3\). Scaling gives equality for
every equilateral triangle. To refute the universal
claim, consider instead the stated triangle with base length
\(2\), height \(h=1/4\), and equal side lengths
\(L=\sqrt{1+h^2}=\sqrt{17}/4\). Its area is \(h\), and
its three altitudes are \(h,2h/L,2h/L\). Since \(L<2\),
the longest is \(H=2h/L=2/\sqrt{17}\).

By symmetry the circumcenter has coordinates \((0,k)\).
Equality of its squared distances to \((1,0)\) and \((0,h)\)
requires \(1+k^2=(h-k)^2\), or
\(k=(h^2-1)/(2h)=-15/8\). Its distance to a base vertex is
\(\sqrt{1+225/64}=17/8\), proving the stated circumradius.

The incenter is also on the vertical symmetry axis. If its
distance to the base is \(r\), its coordinates are \((0,r)\).
Its distance to the side line through \((-1,0)\) and \((0,h)\)
is \((h-r)/\sqrt{1+h^2}\) for the interior point.
Equating this to \(r\) gives
\(r=h/(1+L)=1/(4+\sqrt{17})\), a positive number less than
\(h\). Hence it indeed lies in the interior and touches all
three sides at that distance. Finally, \(R=17/8>2\), while
\(H=2/\sqrt{17}<1/2\). Therefore \(r+R>H\).
Every quantity belongs to an ordinary positive-area triangle;
no infinite radius or degenerate triangle is used as a counterexample.

For the family with equal sides one, choose base endpoints
\((\pm\sin\theta,0)\) and apex \((0,\cos\theta)\).
Its area is \(\sin\theta\cos\theta\), its semiperimeter
is \(1+\sin\theta\), and its altitudes are
\(\cos\theta,2\sin\theta\cos\theta,2\sin\theta\cos\theta\).
Dividing area by semiperimeter gives the inradius, since the
three triangles from the incenter have respective areas equal
to one half their side length times the same inradius.
The circumcenter is \((0,k)\), where equality of squared
distances requires
\(\sin^2\theta+k^2=(\cos\theta-k)^2\).
Solving and taking the distance to the apex gives
\(R=1/(2\cos\theta)\). All denominators are positive in
the stated interval. Continuity of sine and cosine and their
endpoint values give both limits. These limits suggest a
direction in which to search; the preceding concrete triangle
certifies the false universal statement without leaving its domain.
\end{proof}
\end{annexitem}

\begin{annexitem}{annex:p12.heuristic.ships}{thm:p12-heuristic-ships}{Relative motion and the admissible time interval}
\begin{proof}
Subtracting the two positions gives
\((a+tu)-(b+tv)=d+tw\). Subtracting \(b+tv\) from the
whole configuration sends the second position to the origin
and the first to this same relative vector, so all pairwise
distances under consideration are unchanged. If \(w=0\),
the relative vector is constant and every admissible time
has the same distance.

For \(w\ne0\), minimizing the nonnegative distance is equivalent
to minimizing its square. Completing the square gives
\[
 \|d+tw\|^2
 =\|w\|^2(t-t_*)^2+
       \|d\|^2-\frac{(d\cdot w)^2}{\|w\|^2},
 \qquad t_*=-\frac{d\cdot w}{\|w\|^2}.
\]
The coefficient of the squared term is positive, so its unique
minimum over \(\mathbb R\) is at \(t_*\). On the restricted
ray \([0,\infty)\), the allowed point nearest to \(t_*\)
is \(t_*\) if it is nonnegative, and zero otherwise.
This proves the future-time formula, including the case of
a collision when the minimum distance is zero.

The time restriction matters. Let \(a=(0,0)\), \(b=(1,0)\),
\(u=(-1,0)\), \(v=(0,0)\). Then \(d=w=(-1,0)\)
and \(t_*=-1\). The unrestricted trajectories meet at that
past time, giving distance zero. For \(t\ge0\) their distance
is \(1+t\), with minimum one at the initial time. Applying the
unrestricted perpendicular-distance construction without its
time-domain condition would answer a different problem.
\end{proof}
\end{annexitem}
